\begin{solapartat}
Si analitzem la tensió de les bobines del primari i del secundari tindrem:
\[ \left.\begin{aligned} V_p &= N_p\,\frac{\mathrm{d}\Phi}{\mathrm{d}t}\\ V_s &= N_s\,\frac{\mathrm{d}\Phi}{\mathrm{d}t}\end{aligned}\right\} \ \text{\punts{0,3}} \Rightarrow \frac{V_p}{N_p} = \frac{V_s}{N_s} \ \text{\punts{0,4}} \Rightarrow \]
\[ V_s = 4,4\ \mathrm{V} \ \text{\punts{0,3}} \]
\end{solapartat}

\begin{solapartat}
La relació de potències la podem escriure com:
\[ P = P_p = I_pV_p = I_sV_s \ \text{\punts{0,4}} \Rightarrow \]
\[ I_p = I_s\,\frac{V_s}{V_p} = 2\,\mathrm{A} \ \text{\punts{0,3}} \]
\[ P = I_pV_p = 2\,\mathrm{A}\times 220\ \mathrm{V} = 440\ \mathrm{W} \ \text{\punts{0,3}} \]
\end{solapartat}
